\clearpage
\subsection{Reduction to the certificate}
\label{subsec:hko-local-maximum-decision-problem}

The proof is local, and the first step is not computational. The following
elementary reduction is the bridge from the Hausdorff theorem statement to the
dual-vertex coordinates used by the certificate. It is needed because the
exact certificate lives in labelled coordinates, whereas
Theorem~\ref{thm:hko-ten-facet-local-maximum} is a statement about unlabelled
nearby ten-facet polytopes.

\begin{lemma}[Hausdorff-local ten-facet chart]
\label{lem:hko-hausdorff-local-ten-facet-chart}
After shrinking the Hausdorff neighborhood of \(K_{\mathrm{HKO}}\) in
\(\mathcal P_{10}\), every nearby polytope has a unique numbering of its ten
polar vertices by proximity to the fixed ordered HKO polar vertices used by the
certificate:
\[
  K(a)=
  \{x\in\mathbb R^4:\langle a_i,x\rangle\le 1,\ i=1,\ldots,10\},
  \qquad
  a=(a_1,\ldots,a_{10})\in(\mathbb R^4)^{10}.
\]
The supporting inequalities stay facet-defining, boundedness and nonempty
interior persist, and the HKO face combinatorics is unchanged. Conversely, a
sufficiently small coordinate neighborhood of the fixed ordered HKO dual-vertex
list parametrizes this Hausdorff-local stratum.
\end{lemma}

\begin{proof}
The centered HKO representative contains the origin in its interior. After
shrinking the Hausdorff neighborhood, the origin remains in the interior of
every nearby polytope, and polarity is continuous on this neighborhood. Thus a
polytope in \(\mathcal P_{10}\) has ten facets if and only if its polar has ten
vertices.

Since \(K_{\mathrm{HKO}}\) is simple, its polar
\(K_{\mathrm{HKO}}^\circ\) is a simplicial polytope with ten vertices. Choose
pairwise disjoint small neighborhoods of these ten vertices, indexed by the
fixed HKO labelling used in the certificate. For each vertex, choose a linear
functional in the interior of its normal cone. If a ten-vertex polytope \(L\) is
sufficiently Hausdorff-close to \(K_{\mathrm{HKO}}^\circ\), then the maximizer
of the corresponding functional on \(L\) is a vertex in the chosen neighborhood.
Since \(L\) has exactly ten vertices, each neighborhood contains exactly one
vertex and there are no other vertices. Applying this to \(L=K^\circ\), every
Hausdorff-near \(K\in\mathcal P_{10}\) is represented in the fixed labelled
ten-dual-vertex chart.

It remains only to shrink to the stable simple-polytope chart. We use the
standard local stability of simple labelled \(H\)-representations: after
shrinking the labelled coefficients, a simple polytope keeps the same face
lattice. In the present chart this is witnessed by finitely many open
conditions. At \(K_{\mathrm{HKO}}\), each vertex is the transverse intersection
of its four incident supporting hyperplanes, every nonincident inequality has
strict slack at that vertex, the labelled inequalities are irredundant, and
each nonface four-tuple of supporting hyperplanes is infeasible, so no new
vertex can appear from a nonface four-tuple. These determinant, slack,
irredundancy, nonface-infeasibility, boundedness, and origin-interior
conditions remain true for all sufficiently small labelled perturbations, so the
HKO face combinatorics is unchanged in the labelled chart. Conversely, small
perturbations of the fixed ordered HKO polar-vertex list keeping these same open
conditions give the local stratum after shrinking once more.
\end{proof}

This local stability statement is separate from the SageMath certificate: the
certificate begins only after the Hausdorff-local problem has been put in the
fixed labelled dual-vertex coordinates used by the verifier. Hence, by
Lemma~\ref{lem:hko-hausdorff-local-ten-facet-chart}, the coordinate-local
conclusion in this chart proves the stated local result in \(\mathcal P_{10}\).

In a labelled chart, the continuous \(G_{\operatorname{sys}}\)-action induces
a smooth local action on the dual-vertex coordinates. A translation by
\(y\in\mathbb R^4\) sends
\[
  a_i\mapsto \frac{a_i}{1+\langle a_i,y\rangle},
\]
positive scaling by \(\rho>0\) sends \(a_i\mapsto \rho^{-1}a_i\), and a
linear symplectic map \(M\) sends \(a_i\mapsto M^{-T}a_i\). Differentiating
the four translation directions, the scaling direction, and the ten directions
of \(\mathfrak{sp}(4,\mathbb R)\) gives a \(15\)-dimensional orbit tangent
space \(T_{\mathrm{orb}}\subset(\mathbb R^4)^{10}\).

Choose a linear complement \(S\) to \(T_{\mathrm{orb}}\). The exact verifier
checks that \(\dim T_{\mathrm{orb}}=15\), so \(\dim S=25\). After shrinking the
neighborhoods, the map
\[
  (g,s)\longmapsto g\cdot(a_{\mathrm{HKO}}+s)
\]
from a neighborhood of the identity in
\(\mathbb R^4\rtimes(\mathbb R_{>0}\times
\operatorname{Sp}(4,\mathbb R)^0)\) and a
neighborhood of \(0\) in \(S\) is a local diffeomorphism onto a neighborhood of
\(a_{\mathrm{HKO}}\) in the chart. After shrinking, each nearby point has a
unique representation \(g\cdot(a_{\mathrm{HKO}}+s)\) with \(g\) near the
identity and \(s\in S\) small. Thus it is enough to prove that
\(\operatorname{sys}\) has a strict local maximum on \(a_{\mathrm{HKO}}+S\). If
equality holds near HKO, the slice component must be \(s=0\), and the polytope
lies in the \(G_{\operatorname{sys}}\)-orbit. The converse inclusion in the
theorem follows from
Lemma~\ref{lem:preliminaries-systolic-ratio-invariance}.
